\documentclass[12pt]{report}

%------------------------------------------------
% Packages
%------------------------------------------------
\usepackage[a4paper,
left=1.5in,
right=1in,
top=1.5in,
bottom=1in]{geometry}

\usepackage{times}
\usepackage{setspace}
\usepackage{graphicx}
\usepackage{tocloft}
\usepackage{fancyhdr}
\usepackage{titlesec}
\usepackage[hidelinks]{hyperref}

%------------------------------------------------
% Line Spacing
%------------------------------------------------
\onehalfspacing

%------------------------------------------------
% Header and Footer
%------------------------------------------------
\pagestyle{fancy}
\fancyhf{}
\fancyfoot[C]{\thepage}
\renewcommand{\headrulewidth}{0pt}

%------------------------------------------------
% Chapter Formatting
%------------------------------------------------
\titleformat{\chapter}
{\normalfont\huge\bfseries}
{\thechapter}{20pt}{\Huge}

%------------------------------------------------
% Begin Document
%------------------------------------------------
\begin{document}

%------------------------------------------------
% Title Page
%------------------------------------------------
\begin{titlepage}
    \centering
    
    \vspace*{1.5cm}
    
    {\Huge \textbf{Field Insights Report}}\\[1.5cm]
    
    \vspace{2cm}

    \vspace*{1.0cm}
    
    { \textbf{A Better Open Area}}\\[1.5cm]
    
    \vspace{1cm}
    
    {\Large
    \textbf{How the UCSC Open Area is used for academic and non-academic activities across different times of the day, and what factors influence students to choose the space for studying versus social interaction}
    }\\[2cm]
    
    \vfill
    
    \begin{flushleft}
    \textbf{Research Area:} UCSC Open Area\\[0.5cm]
    
    \textbf{Module:} SCS 3312 -- Research Methods\\[0.5cm]

    \textbf{Group Number:} 08\\[0.5cm]
    
    
    \textbf{Individual Report}\\
    Member name:  K.B Jayawardhana - 2023 CS 078\\[0.5cm]
    
    \textbf{University:} University of Colombo School of Computing\\[0.5cm]
    
    \textbf{Date:} 11.06.2026
    \end{flushleft}
    
\end{titlepage}

%------------------------------------------------
% Roman Numbering
%------------------------------------------------
\pagenumbering{roman}

%------------------------------------------------
% Acknowledgement
%------------------------------------------------
\chapter*{Acknowledgement}
\addcontentsline{toc}{chapter}{Acknowledgement}

I am thankful to the students of the University of Colombo School of Computing for participating in this study by giving up their time and views in the interviews. Their valuable inputs helped validate 
the findings obtained from the observations made during the research study.

I would like to thank the module coordinator of the module SCS 3312, Research Methods for all his support through out the research process and my group members for their collaboration during the field work.
\newpage

%------------------------------------------------
% Table of Contents
%------------------------------------------------
\tableofcontents
\newpage

%------------------------------------------------
% List of Figures
%------------------------------------------------
%\listoffigures
%\newpage

%------------------------------------------------
% List of Tables
%------------------------------------------------
%\listoftables
%\newpage

%------------------------------------------------
% Abstract Chapter
%------------------------------------------------
\chapter*{Abstract}
\addcontentsline{toc}{chapter}{Abstract}

This report highlights the results of a field study undertaken at the Open Area of the University of Colombo School of Computing (UCSC). The research focuses on how the space is used for both academic and non-academic purposes at various hours of the day, and what determines the choice of the space either for studying or engaging in socializing. A qualitative research design was applied, including 11 sessions of non-participant observations and 11 semi-structured interviews. Major findings revealed that there was a greater tendency 
for academic pursuits at the morning and evening hours while the midday and afternoon time were marked by non-academic activities. Heat, noise, lightening, and other environmental variables were determined as critical behavioral determinants among others.

\newpage


%------------------------------------------------
% List of Acronyms
%------------------------------------------------
\chapter*{List of Acronyms}
\addcontentsline{toc}{chapter}{List of Acronyms}

\begin{tabular}{ll}
\textbf{UCSC} & University of Colombo School of Computing \\
\end{tabular}

\newpage

%------------------------------------------------
% Arabic Numbering Starts
%------------------------------------------------
\pagenumbering{arabic}

%------------------------------------------------
% Introduction Chapter
%------------------------------------------------
\chapter{Introduction}

The UCSC Open Area is regularly used by many students during the course of the day. Nevertheless, the behavior of these students while using this area greatly differs from each other. While some of the students use this area to study individually or in groups, some other students find this area ideal for engaging in social conversation, club activities, or even for recreation. Unfortunately, very little is known about why students change their behaviors like this. It is hard to assess the efficacy of the use of such an area without understanding the reasons behind student behavior.

Open areas serve an important role in university life as places of informal education where students interact with each other academically and socially. Understanding this specific area and the student behavior patterns which influence its use may help in fitting such an area to the performed activity.

\vspace{0.5cm}
\textbf{Field Site:}

The selected field site for this study is the Open Area of the University of Colombo School of Computing. Observations were conducted during multiple periods of the day (within operating hours of the university), including:

\begin{itemize}
    \item Morning
    \item Midday
    \item Afternoon
    \item Evening
\end{itemize}

\vspace{0.5cm}
\textbf{Research Questions:}

\begin{enumerate}
    \item How does the utilization of the UCSC Open Area vary across different times of the day?
    \item What proportion of observed activities within the open area are academic versus non-academic?
    \item What environmental and social factors influence students' decision to use the space for studying rather than social interaction?
\end{enumerate}

\vspace{0.5cm}
\textbf{Methodology:}

This phase was performed using the qualitative approach, whereby the primary emphasis was laid on the analysis of behaviors, interactions, and perceptions rather than the statistical measurement of relationships. Several observations were carried out by different group members in the UCSC Open Area. The semi-structured interviews were conducted on willing participants depending on the availability during observations.
\begin{itemize}
    \item Total number of observation sessions conducted: 11
    \item Total number of interviews conducted: 11
\end{itemize}

Observation notes were documented immediately during fieldwork sessions and interview responses were recorded in written form. In addition, environmental and contextual data about the observed space were documented alongside behavioural observations.

\chapter{Observation summaries}

Following are summaries of selected observation sessions.

\section{Observation Session 1 -- Midday, 26 May 2026}
\textbf{Conditions: Sunny, warm, high noise, moderate lighting, medium occupancy}\\
The space was densely occupied. The loud and active assembly of individuals who were building lanterns for the Vesak celebration near the tree area caused noise and disturbance. The students involved in study activities chose to sit at tables that were located farther away from the event.
A group of students enguaged in university club organizational work was observed and confirmed through an interview.

\section{Observation Session 2 -- Morning, 2 June 2026}
\textbf{Conditions: Cloudy, moderate temperature, low noise, Dim lighting, low occupancy}\\
The activity in the room was mainly academic-related, with students seen reviewing notes and reading their course materials. Many took positions at places that had charging stations. It seemed like the quiet and less-crowded atmosphere made for an environment suitable for studying.

\section{Observation Session 3 -- Midday, 2 June 2026}
\textbf{Conditions: Sunny, warm, moderate noise, moderate lighing, medium occupancy}\\
A blend of academic and non-academic uses was achieved. This served as an intermediate zone for lectures. Students using academics were concentrated around areas where they could charge their devices, whereas those involved in social activities used the open spaces.

\section{Observation Session 4 -- Afternoon, 5 June 2026}
\textbf{Conditions: Cloudy, warm, moderate noise, dim lighting, moderate occupancy}\\
The academic crowds gravitated towards the science faculty end, whereas the social groups remained at the end facing UCSC. Those with textbooks left as the lights faded out. Some moved to areas where there was enough light.

\chapter{Interview summaries}

\section{Reasons for Choosing the Space}
It was found that students perceived this open space as tranquil and easily accessible when compared to other spaces on campus. As one student mentioned, "this is a very tranquil place compared to other places for studying." Others pointed out "trees and greenery around". The presence of friends was also a recurring reason for staying, with several participants stating they remained because they were working or socialising with others.


\section{Physical Facilities}
Device charging emerged as the most frequently mentioned physical determinant affecting room selection for students. Room seating could also accommodate individual users as well as those working in groups.

\section{Environmental Influences}
Heat was the most frequently mentioned. Some of the students complained that, they left the place at some point during warmer times. Noise associated with social or extracurricular activities is yet another factor mentioned by those who use the place for study.

\section{Areas for Improvement}
Fan installation emerged as the most frequent recommendation, made by the participants irrespective of their activities. Lighting needed to be turned on sooner and more consistently in the evening. While other aesthetic enhancements such as adding plants were recommended, the main concern remained ventilation.

\chapter{Findings}
This section focuses on the information derived from the observation sessions and the interviews.

\section{Types of activities observed}
Multiple observations across different sessions have shown that the UCSC Open Area is a venue for diverse academic and non-academic activities. The academic activities carried out in this area include studying individually and in groups, 
revision of lessons through book reading, note-taking and laptop use, and academic discussions. Some students used this space to complete small academic submissions. However most students engaged in academic activities did so for a relatively 
short period of time(less than 2 hours) with some students changing from studying to mobile phone use, sleeping/resting. The non-academic activities undertaken include socializing with friends, doing club/organization work, preparing for events, 
leisure time activities, waiting or resting during lecture intervals.

It was recorded that a combination of academic and non-academic activities occurred in all of the observation sessions. While some students were engaged in focused academic work, others were simultaneously involved in social discussions, extracurricular 
planning, or recreational activities. However academic activities were dominant in the morning and the evening sessions while non-academic activities were dominant in the mid-day and Afternoon sessions. The occupancy level of the space was at a moderate 
to high in the mid day and evening sessions while it was low in the morning and evening sessions.

\vspace{0.9cm}

\section{Spatial usage}
Students did not use the space evenly. Students involved in academic work frequently selected the side towards the science faculty. During several sessions, study groups and individual learners were observed occupying locations away from major sources 
of noise and pedestrian movement. Areas with easier access to charging points also appeared to attract students engaged in academic activities. In contrast, the socially engaged students and those who participated in extracurricular activities usually
located themselves within more centralized areas. During one of the observation sessions, there was a group of students constructing Vesak lanterns at several tables by the tree-area towards the UCSC. This became the focus of social activity, whereas 
academic users located themselves further away from this activity.

\section{Environmental conditions}
Temperature was one of the main issues for the students. The observations support this as most of the sessions recorded moderate to high temperatures. From the interviews, it can be derived that adding fans to the space might improve students decision 
in choosing the space for academic and non-academic activities and help retain them for a longer time in the space.

Lighting condition in the evening are inadequate. As the observations and the interviews suggest the lights in the space should be turned in earlier than the current time and consistently.

Noise levels also affected behavior. There was increased noise when sessions were held that involved any group activity, events planning, or socializing. At such times, students who were studying could be seen sitting in quieter parts of the area. 
However, during quiet times, there seemed to be an even distribution of the academic activity in the area. This leads to the assumption that non-academic activities hinder the academic activities within the space and vice-versa.

Weather conditions also affected the usability of the space. Rainy conditions reduced overall activity levels in some sessions and increased the number of individuals waiting or resting in sheltered areas. Several observations indicated that weather 
exposure remains a limitation of the open area despite its popularity among students. This might be because of the inadequate drainage system in the space.



%------------------------------------------------
% Critical Analysis
%------------------------------------------------
\chapter{Critical Analysis}

Several limitations must be acknowledged when interpreting these findings of these observations.

The method of observation used in the research involved mostly non-participant observation and semi-structured interviewing techniques. Both methods are associated with observer bias and social desirability bias, respectively, because students who were under observation might have changed their behavior, and the participants who underwent interviews might have produced responses that would sound better and be more socially desirable than truthful.
Also, no observation or interview sessions took place at the weekends.

Despite such shortcomings, the convergence of results obtained during both the observations and interviews serves to bolster the reliability of the major conclusions. Themes of environmental discomfort, preference for quiet areas by academics, and the versatility of the area emerged repeatedly throughout multiple sessions and through multiple subjects.



\chapter{Reflection}

Before undertaking the field work, it was thought that the UCSC Open Area was utilized by students mainly for academics and social engagements. Nonetheless, through the field work, it became clear that these two main types of actions entail many other different forms of actions. It is therefore clear that students use this area to live together in dynamic ways.

One of the most important realizations while conducting the field work was that students' actions cannot always be easily comprehended solely through observation. During a number of observations, it was very challenging to differentiate between discussions about academics, extra-curricular activities, and general chatting based only on observation. The interview 
sessions came in handy in such instances.

This study was limited to qualitative data collection. Future work should incorporate the quantitative phase outlined in the group research proposal, including structured headcounts, activity frequency tallies, and environmental measurements such as temperature and noise levels recorded across sessions. This would allow for statistical validation of the patterns identified here.


%------------------------------------------------
% References
%------------------------------------------------
%------------------------------------------------
% Index
%------------------------------------------------
%------------------------------------------------
% End Document
%------------------------------------------------
\end{document}
